\section{Introduction}
\label{sec:introduction}

Product formulas remain among the most transparent constructions for digital
Hamiltonian simulation.  If $H=\sum_{\gamma=1}^{\Gamma}H_\gamma$, a product
formula replaces $\exp(-itH)$ by a sequence of exponentials of simpler
fragments $H_\gamma$.  The method needs no block encoding, can preserve the
geometry of local interactions, and maps naturally to hardware with limited
connectivity.  Its practical cost, however, is determined not merely by its
formal order but by the constant in a rigorous finite-time error bound.  A
factor lost in this constant is amplified when the bound is inverted to select
an integer number of Trotter steps and then multiplied by every local
propagator in every step.

The modern commutator-scaling theory of Childs, Su, Tran, Wiebe, and Zhu
provides a general foundation for such estimates
\cite{childs2021trotter}.  It replaces purely norm-sum scaling by sums of
nested commutators and thereby exploits locality and exact commutation.  This
framework is sufficiently general to cover arbitrary order and many
Hamiltonian classes, but a general theorem must remain conservative about the
concrete algebra of a fixed model.  In a Pauli lattice, multiple formal
commutator paths may evaluate to the same operator with opposite signs;
translation-equivalent terms may repeat; and distinct surviving Paulis may
anticommute, so their sum has a Euclidean rather than taxicab norm.  Applying a
triangle inequality before these facts become visible irreversibly discards
them.

This paper studies the resulting ``last-mile'' problem:

\begin{quote}
Given a fixed local Hamiltonian, fragmentation, product formula, simulation
time, error tolerance, and compilation model, how can one automatically
produce a tighter \emph{machine-checkable} worst-case resource bound without
changing the target quantum circuit family?
\end{quote}

Our answer is to treat error analysis as compilation.  The source program is
the Hamiltonian and product formula.  The compiler lowers it through
successive algebraic intermediate representations, performs semantics-
preserving optimization passes, and emits a proof object.  In contrast with a
normal circuit optimizer, the principal output is not a different sequence of
quantum gates.  It is a smaller certified number of repetitions of the same
sequence.  The design follows a proof-carrying principle: expensive symbolic
search may be heuristic, but an accepted result must carry all data required
for a simpler independent verifier to reconstruct the logical obligations.

\subsection{Contributions}

The main contributions are as follows.

\begin{enumerate}
\item \textbf{A multi-representation error compiler.}  We compute the
  truncated logarithm of the complete product formula in a sparse free
  associative algebra, project its homogeneous defects to free-Lie elements,
  evaluate those elements in a concrete symplectic Pauli representation, and
  finally emit exact rational certificate data.  Each lowering step has an
  explicit invariant and can be checked independently.

\item \textbf{Norm-last concrete evaluation.}  We canonicalize finite Pauli
  configurations under lattice translation and combine identical Pauli
  strings with exact coefficients before applying a final norm inequality.
  This reverses the conventional expand-and-bound order and preserves
  cancellations that are invisible at the level of abstract support counts.

\item \textbf{Certified anticommuting norm compression.}  We construct a
  graph whose vertices are leading-defect Pauli terms and whose edges encode
  anticommutation.  An untrusted search partitions terms into pairwise-
  anticommuting groups.  The certificate records the full partition and
  outward rational square-root bounds; the verifier rechecks exact coverage
  and every edge used.  For the principal instance, this lowers the D4
  translation-cell bound from 20.160968407335066 to 6.472926505087888.

\item \textbf{A nonasymptotic right-generator closure.}  We convert the
  degree-five and degree-seven logarithmic defects into the right logarithmic
  generator, retain degrees D4 through D7, use a Heisenberg-specific local
  commutator lemma, and control all remaining degrees by a rational geometric
  locality tail.  Thus the resource result does not rely on leading-order
  extrapolation.

\item \textbf{An independently reproducible finite-instance result.}  For a
  $12\times12$ periodic Heisenberg lattice, the certificate accepts 97 steps
  and rejects 96.  Relative to a pinned 393-step instantiation of the
  published high-order theorem, this reduces merged group exponentials from
  11,791 to 2,911.  A fast verifier, deep regeneration mode, small-system
  dense cross-check, mutation tests, and a SHA-256 manifest support the claim.
\end{enumerate}

\subsection{Scope and nonclaims}

The result is a finite-size constant-factor improvement, not a new asymptotic
Hamiltonian-simulation complexity.  It applies to a uniform operator-norm
guarantee for a fixed deterministic product formula.  Methods based on low-
energy inputs, average-case norms, multi-product formulas, linear combinations
of unitaries, extrapolation, or stochastic compensation solve different
resource problems and may achieve stronger precision scaling under their own
cost models.  We compare these directions in \cref{sec:related} but do not
rank incomparable gate counts as a single global record.

The exact D5 side study and the $r=78$ arithmetic indicate a possible
fivefold threshold.  They do not establish it.  Throughout this manuscript,
``certified'' refers only to statements accepted by the supplied rational
certificate and verifier.  Numerical dense simulation is used solely as a
cross-check, and conditional research estimates are labeled explicitly.
